\newpage{\ }
\cleardoublepage

\chapter{Aplicación Móvil}
\label{ch:aplicacion-movil}

Para conectar los motores de generación de oraciones con las personas que los necesitan nace la aplicación móvil Dixi.

El diseño de la interfaz incorpora decisiones de accesibilidad condicionadas por el perfil de los usuarios de CAA: personas con Trastorno del Espectro Autista, parálisis cerebral, afasia, ELA u otras condiciones que dificultan el habla. La app ofrece retroalimentación visual y háptica, tamaños ajustables, contraste adecuado y una organización de pictogramas pensada para construir mensajes con el menor número posible de toques.

El sistema sigue una arquitectura cliente-servidor:

\begin{itemize}
    \item El cliente es la aplicación móvil, desarrollada con React Native~\cite{reactnative2024} y Expo~\cite{expo2024}. Presenta el tablero de pictogramas, recoge la selección del usuario y muestra la frase generada.
    \item El servidor es una API REST desarrollada con FastAPI~\cite{fastapi2024}. Recibe las etiquetas de los pictogramas seleccionados, las pasa al motor de generación correspondiente y devuelve la frase resultante.
\end{itemize}

Los motores de generación están escritos en Python y no pueden ejecutarse en un dispositivo móvil, lo que justifica esta separación. Las secciones siguientes detallan las tecnologías elegidas para cada lado de esta arquitectura y cómo se integran entre sí.

\section{Decisiones tecnológicas}
\label{sec:app-tecnologias}

La elección de React Native~\cite{reactnative2024} como framework de la app responde a la necesidad de cubrir iOS y Android desde un único código fuente. El proyecto usa React Native 0.81 con React 19, compilado mediante el SDK de Expo 54~\cite{expo2024}. Expo aporta un conjunto de bibliotecas nativas preconfiguradas (cámara, audio, hápticos, almacenamiento) y un sistema de compilación que abstrae la configuración de Xcode y Gradle.

La navegación se gestiona con Expo Router 6, que organiza las pantallas como ficheros en un directorio, de forma análoga a Next.js en el entorno web. La app tiene habilitada la \textit{New Architecture} de React Native, que emplea el motor Hermes para la ejecución de JavaScript. Todo el código del cliente está escrito en TypeScript 5.9, lo que añade tipado estático sobre JavaScript.

El servidor usa FastAPI~\cite{fastapi2024}, un framework web para Python. La elección viene condicionada por los motores de generación: los tres están escritos en Python y dependen de bibliotecas que no existen en otros lenguajes. Portar cualquiera de estos componentes a TypeScript o Go habría supuesto reescribirlos desde cero, y alternativas como AWS Lambda quedaban descartadas porque el modelo de embeddings excede los límites de memoria de las funciones \textit{serverless}.

La decisión de interponer un servidor entre la app y los motores tiene una motivación adicional de seguridad. El motor de LLMs requiere una clave de API de OpenAI. Si la app llamase directamente a la API de OpenAI, esa clave tendría que estar embebida en el binario de la aplicación (APK en Android, IPA en iOS), donde sería extraíble mediante ingeniería inversa. Con el servidor como intermediario, la clave reside exclusivamente en el entorno del servidor y la app nunca la conoce.

El frontend se apoya en bibliotecas nativas de Expo y React Native. Las imágenes de pictogramas se cargan con \texttt{expo-image}, que cachea en disco y recicla vistas para contener el consumo de memoria cuando el usuario desplaza la cuadrícula. Las animaciones de retroalimentación táctil corren en el hilo nativo gracias a \texttt{react-native-reanimated} 4.1, manteniendo 60 fotogramas por segundo incluso en dispositivos de gama baja.

Para la síntesis de voz la app usa \texttt{expo-speech} 14.0, que da acceso a los motores nativos de cada plataforma (AVSpeechSynthesizer en iOS, TextToSpeech en Android). % TODO: ampliar cuando se conecte con TTS externo (ElevenLabs u otro)
La retroalimentación háptica se gestiona con \texttt{expo-haptics} 15.0, y las listas de pictogramas con \texttt{@shopify/flash-list} 2.0, que mejora el rendimiento de desplazamiento respecto al FlatList estándar de React Native.

\section{Arquitectura del sistema}
\label{sec:app-arquitectura}

El flujo de comunicación completo atraviesa cuatro etapas:

\begin{enumerate}
    \item El usuario selecciona pictogramas en el tablero de la app.
    \item La app extrae las etiquetas textuales de los pictogramas seleccionados y las envía al servidor como una lista de cadenas de texto.
    \item El servidor delega la generación al motor configurado y devuelve la frase resultante.
    \item La app muestra la frase y, si el usuario lo solicita, la reproduce como voz.
\end{enumerate}

% TODO: sustituir el ejemplo textual por un diagrama de flujo visual
% TODO: actualizar ejemplo cuando se elija la técnica definitiva
Por ejemplo, si el usuario toca los pictogramas ``yo'', ``querer'' y ``agua'', la app los envía al servidor, que genera la oración ``Yo quiero agua'' y la devuelve junto con el tiempo de procesamiento y un indicador de éxito. La app muestra la frase en la barra de construcción.

El servidor organiza los motores de generación mediante dos patrones de diseño complementarios: \textit{Strategy}~\cite{refactoringguru_strategy} y \textit{Registry}~\cite{refactoringguru_registry}. Todos los motores implementan una misma interfaz abstracta que define tres operaciones: cargar los recursos necesarios al arrancar (modelos, conexiones a APIs externas), generar una oración a partir de una lista de pictogramas, e identificarse con un nombre único. La respuesta de cada motor incluye la oración generada, el tiempo de procesamiento, un indicador de éxito y, opcionalmente, un mensaje de error o metadatos adicionales.

Un registro central mantiene la relación entre nombres de motores y sus implementaciones. Al arrancar, el servidor lee de sus variables de entorno la lista de motores que debe cargar, los instancia y ejecuta su carga de recursos. Cada motor se inicializa una sola vez y atiende todas las peticiones posteriores como \textit{singleton}~\cite{refactoringguru_registry}. Añadir un motor nuevo consiste en implementar la interfaz abstracta, registrarlo e incluir su nombre en la configuración, sin modificar el código existente.

% TODO: reescribir este párrafo cuando se integre ElevenLabs u otro TTS externo.
% Entonces la síntesis principal será vía servidor y el motor nativo pasará a ser fallback offline.
La última etapa del flujo es la síntesis de voz. Actualmente la app utiliza los motores nativos de cada plataforma (iOS y Android) para reproducir la frase generada como audio. Esta solución es provisional: la voz definitiva procederá de un servicio de síntesis neuronal (ElevenLabs u otro similar) que produzca locuciones más naturales, y al que la app accederá a través de una ruta del servidor. Los motores nativos quedarán entonces como fallback para situaciones sin conexión a Internet. La pantalla de ajustes ya permite seleccionar voz, velocidad y tono, parámetros que se mantendrán independientemente del motor subyacente.

Con esta arquitectura definida, las secciones siguientes describen los detalles concretos de cada lado: primero el servidor y después la app.

\section{Servidor de generación}
\label{sec:app-servidor}

El servidor se organiza en tres módulos: uno con el punto de entrada, la configuración y los esquemas de datos; otro con la interfaz abstracta, el registro y los seis adaptadores de motores; y un tercero con las rutas HTTP. La configuración se lee de variables de entorno con valores por defecto. Durante el desarrollo se usa el motor de gramática formal como motor por defecto, ya que no requiere dependencias externas (ni modelo de embeddings ni clave de API) y permite probar el flujo completo de la aplicación sin infraestructura adicional. % TODO: actualizar cuando se elija la técnica definitiva a partir de los resultados de la evaluación comparativa (sección~\ref{ch:evaluacion-comparativa})

El endpoint principal recibe una lista de pictogramas (entre 1 y 20 etiquetas) y un parámetro que indica qué motor usar. % TODO: indicar aquí el motor por defecto cuando se elija la técnica definitiva La generación se ejecuta en un \textit{thread pool} con un \textit{timeout} configurable de 30 segundos: si el motor solicitado no está disponible el servidor responde con un error 400, y si la generación excede el tiempo límite, con un error 408. El servidor expone además dos endpoints de salud, uno que indica si el proceso está vivo y otro que verifica si al menos un motor está operativo.

Los recursos pesados (modelos de embeddings, gramáticas, modelos de perplejidad) se cargan una sola vez al arrancar el servidor y permanecen en memoria durante toda la vida del proceso, compartiéndose entre peticiones. Sin este mecanismo, cada petición tendría que recargar modelos de varios gigabytes, lo que haría inviable el servicio.

La clave de API de OpenAI reside exclusivamente en la configuración del servidor. La app móvil no la contiene ni la necesita: envía las etiquetas de pictogramas y recibe la frase generada, sin conocer qué motor la produjo ni qué credenciales se usaron. El servidor no implementa autenticación propia, una simplificación aceptable en un proyecto académico donde se ejecuta en red local.

El manejo de errores sigue una estrategia de degradación. Si un motor falla durante la generación, el adaptador devuelve la respuesta con un indicador de error en lugar de provocar un fallo completo del servidor. Esto permite que la app distinga entre un fallo de red (que requiere reintentar) y un fallo del motor (que puede resolverse mostrando un mensaje al usuario o recurriendo al generador local como respaldo).

El servidor está contenerizado con Docker~\cite{docker2024} para facilitar su despliegue. La imagen parte de una distribución ligera de Python 3.13, copia el código de la API y el de la técnica de generación seleccionada % TODO: actualizar cuando se elija la técnica definitiva, instala las dependencias e inicia el servidor en el puerto 8000. Un fichero de configuración de Docker Compose permite levantar el servicio con un solo comando, sin necesidad de instalar Python ni las dependencias en la máquina anfitriona. La contenerización convierte el despliegue en una operación reproducible: la misma imagen que se ejecuta en la máquina del desarrollador puede subirse a un registro y desplegarse en cualquier servicio \textit{cloud} (AWS ECS, Google Cloud Run, Azure Container Instances) sin ajustes adicionales, ya que el contenedor encapsula el intérprete, las dependencias y el código en un artefacto autocontenido.

\section{Interfaz de usuario}
\label{sec:app-interfaz}

Los usuarios de CAA tienen con frecuencia dificultades motoras o visuales que condicionan cómo interactúan con una pantalla táctil. El diseño de Dixi subordina las decisiones estéticas a la accesibilidad: todas las pantallas incorporan retroalimentación multisensorial configurable (visual, háptica y auditiva), tamaños ajustables y paletas con contraste adecuado. La app organiza sus tres pantallas mediante una barra de pestañas inferior, el patrón de navegación estándar en aplicaciones móviles de uso frecuente.

La pantalla principal es el tablero de comunicación, el espacio donde el usuario construye mensajes seleccionando pictogramas. En la parte superior, una barra de construcción muestra en \textit{scroll} horizontal los pictogramas que el usuario ha ido seleccionando, con la frase generada debajo. Tocar un pictograma en esta barra lo elimina de la selección, permitiendo corregir sin empezar de nuevo. Debajo, un campo de búsqueda permite filtrar el catálogo de pictogramas por nombre.

Una barra de categorías muestra las nueve categorías del catálogo como píldoras horizontales con los colores del sistema de FitzGerald~\cite{fitzgerald1949straight}, un código cromático estándar en CAA que asigna un color a cada categoría gramatical: naranja para personas, verde para acciones, rosa para necesidades, naranja oscuro para objetos, azul para lugares, morado para cualidades, rojo para emociones, marrón para tiempo y gris para preguntas. El cuerpo de la pantalla muestra los pictogramas en una cuadrícula configurable de tres a cinco columnas. Cada tarjeta muestra la imagen del pictograma cargada desde el CDN de ARASAAC~\cite{arasaac} y su etiqueta textual, y al tocarla proporciona retroalimentación visual y opcionalmente háptica. Dos botones flotantes completan la pantalla: uno para añadir pictogramas personalizados mediante un modal y otro para reproducir la frase generada con síntesis de voz. El catálogo incluye más de 200 pictogramas ARASAAC organizados en las nueve categorías, más una categoría ``Todas'' que los muestra sin filtrar. Los colores de FitzGerald se implementan con dos variantes (tema claro y tema oscuro), con valores cromáticos ajustados en cada modo para mantener ratios de contraste adecuados. El código cromático tiene una función práctica: permite al usuario localizar la categoría que busca sin leer las etiquetas, lo que beneficia a personas con dificultades lectoras.

% TODO: incluir figura con captura del tablero de comunicación

La segunda pantalla contiene frases predefinidas para comunicación rápida. El usuario no necesita componer pictogramas: toca una frase y el sistema la reproduce directamente. Las frases se organizan por contexto de uso (casa, colegio, médico, social, emociones) en una lista con secciones colapsables. Cada frase tiene un botón de reproducción y un indicador de favorito. El usuario puede añadir frases propias a una sección personal (``Mis frases'') y buscar entre todas las frases con filtrado en tiempo real. Esta pantalla complementa al tablero de comunicación: el tablero permite construir mensajes nuevos con pictogramas, mientras que las frases predefinidas cubren expresiones de uso frecuente que el usuario quiere tener siempre a mano.

% TODO: incluir figura con captura de la pantalla de frases predefinidas

La tercera pantalla agrupa los ajustes de la aplicación. La sección de voz permite seleccionar una voz del sistema (filtrada a voces en español), ajustar la velocidad de habla entre 0,5x y 2,0x y modificar el tono entre 0,5 y 2,0, con un botón para probar la configuración antes de aplicarla. La sección de apariencia controla el modo claro u oscuro, el tamaño de los pictogramas (pequeño, mediano o grande) y el número de columnas de la cuadrícula. La sección de accesibilidad permite activar o desactivar la vibración háptica y los efectos de sonido, y ajustar el tamaño de fuente entre 14 y 22 píxeles. La pantalla cierra con una sección informativa que muestra el logotipo de Dixi, la versión de la aplicación y los créditos con un enlace a ARASAAC. Todos los elementos interactivos de la app están etiquetados para lectores de pantalla, y la interfaz respeta las áreas seguras del dispositivo para que el contenido no quede oculto por el \textit{notch} o la barra de estado.

% TODO: incluir figura con captura de la pantalla de ajustes

